\renewcommand{\thechapter}{ข}
\chapter{โค้ดต้นแบบ MicroPython และไลบรารี Modbus CRC-16}
\label{chap:appendixB}

\section{ไลบรารี Modbus RTU สำหรับชิป ESP32-S3}

ซอร์สโค้ดภาษา MicroPython สำหรับการสื่อสารกับโพรบวัดดิน 7-in-1 และอินเวอร์เตอร์ขับมอเตอร์ VFD ผ่านพอร์ตซีเรียล UART2

\begin{pythonbox}[โมดูล modbus\_crc16.py สำหรับ MicroPython]
import struct
from machine import UART, Pin

class ModbusRTUMaster:
    def __init__(self, uart_id: int = 2, tx_pin: int = 17, rx_pin: int = 16, 
                 baudrate: int = 9600, de_re_pin: int = 4):
        self.uart = UART(uart_id, baudrate=baudrate, tx=Pin(tx_pin), rx=Pin(rx_pin))
        self.de_re = Pin(de_re_pin, Pin.OUT)
        self.de_re.value(0) # ตั้งค่ารับเริ่มต้น (Receive Mode)

    @staticmethod
    def calculate_crc(data: bytes) -> int:
        crc = 0xFFFF
        for byte in data:
            crc ^= byte
            for _ in range(8):
                if crc & 0x0001:
                    crc = (crc >> 1) ^ 0xA001
                else:
                    crc >>= 1
        return crc

    def send_command(self, slave_id: int, func_code: int, start_reg: int, count: int) -> bytes:
        raw_cmd = struct.pack(">BBHH", slave_id, func_code, start_reg, count)
        crc = self.calculate_crc(raw_cmd)
        full_packet = raw_cmd + struct.pack("<H", crc)
        
        # สลับเป็นโหมดส่ง (Transmit Mode)
        self.de_re.value(1)
        self.uart.write(full_packet)
        # รอให้ส่งข้อมูลเสร็จสิ้นแล้วสลับกลับเป็นโหมดรับ
        self.de_re.value(0)
        
        # อ่านข้อมูลตอบกลับภายใน 1 วินาที
        return self.uart.read(256)
\end{pythonbox}
\clearpage
